% =====================================================================
%  2bat-ud4-15.tex  —  SESSIÓ 15: Taller de preparació de la prova
%  (sense preàmbul: s'inclou des de main.tex)
% =====================================================================
\horatitol{Sessió 15 — Taller de preparació de la prova}%
{Conèixer l'estructura de la prova i la rúbrica, assajar un simulacre per blocs i fer un pla d'estudi personal.}

\subsection{Idea clau}

Ens preparem per a la prova \textbf{sense sorpreses}. La prova de la sessió següent té
\textbf{quatre blocs}, un per cada criteri (C1-C4). Saber exactament què es demana i
com es corregeix fa que l'examen deixi de ser imprevisible. Després assajarem amb un
\textbf{simulacre} amb temps controlat.

\begin{idea}
\textbf{Estructura de la prova} (un bloc per criteri):
\begin{enumerate}[leftmargin=1.6em, itemsep=2pt]
  \item \textbf{Bloc 1 (C1) — Càlcul de derivades:} derivar un polinomi i avaluar la
        derivada en un punt.
  \item \textbf{Bloc 2 (C2) — Creixement:} estudiar on una funció creix i decreix amb el
        signe de $f'$.
  \item \textbf{Bloc 3 (C3) — Màxims i mínims:} trobar i classificar els punts crítics.
  \item \textbf{Bloc 4 (C4) — Optimització:} plantejar i resoldre un problema de context
        social i interpretar-ne la solució.
\end{enumerate}
\textbf{Es valora}, a més del resultat: el \textbf{plantejament} correcte de la funció,
la \textbf{classificació justificada} (signe de $f'$) i la \textbf{interpretació} en
context.
\end{idea}

\subsection{Rúbrica simplificada}

\noindent Així es corregeix cada bloc (nivells d'assoliment):
\begin{center}
\renewcommand{\arraystretch}{1.3}
\begin{tabular}{p{2.4cm} p{10.2cm}}
\toprule
\textbf{Nivell} & \textbf{Què demostra} \\
\midrule
\textbf{AE} (excel·lent) & Resol sense errors i interpreta el resultat en el context
social. \\
\textbf{AN} (notable) & Resol correctament amb algun error menor; la interpretació és
essencialment bona. \\
\textbf{AS} (suficient) & Resol el procediment bàsic amb alguns errors; interpretació
incompleta. \\
\textbf{NA} (no assolit) & No identifica el procediment o comet errors que
n'invaliden el resultat. \\
\bottomrule
\end{tabular}
\end{center}

\subsection{Simulacre — un ítem de cada bloc}

\begin{exemple}[com es resol un ítem ben fet]
\textbf{Ítem tipus (Bloc 3):} troba i classifica els punts crítics de
$f(x)=x^3-12x+5$. Un \textbf{ítem complet} deriva, troba els crítics i \textbf{justifica}
la classificació amb el signe:
\[
f'(x)=3x^2-12=3(x-2)(x+2)=0 \Rightarrow x=-2,\,x=2;
\]
en $x=-2$, $f'$ passa de $+$ a $-$ (màxim $(-2,21)$); en $x=2$, de $-$ a $+$ (mínim
$(2,-11)$). \emph{Classificar amb el canvi de signe, i no només trobar on $f'=0$,
distingeix un AE d'un AS.}
\end{exemple}

\begin{exercicis}
\textbf{Resol el simulacre amb temps controlat. Un ítem per bloc:}
\begin{exlist}
  \item \textbf{(Bloc 1 — C1)} Deriva $f(x)=x^3-5x^2+2x-8$ i calcula $f'(2)$. Digues si,
        en $x=2$, la funció creix o decreix.
  \item \textbf{(Bloc 2 — C2)} Estudia on creix i on decreix $f(x)=x^3-3x^2-9x+2$, amb
        taula de signes.
  \item \textbf{(Bloc 3 — C3)} Troba i classifica els punts crítics de $f(x)=x^3-12x+5$.
  \item \textbf{(Bloc 4 — C4)} Una entitat tanca un recinte rectangular contra una paret
        amb $36$ m de tanca (tres costats). Si els costats perpendiculars fan $x$, troba
        les mides de màxima superfície i interpreta el resultat.
\end{exlist}
\end{exercicis}

\subsection{Formulari-resum (suport)}

\begin{idea}
\textbf{Tingues-ho a mà mentre prepares la prova:}
\begin{itemize}[leftmargin=1.4em, itemsep=2pt]
  \item \textbf{Derivar un polinomi:} $x^n\to n\,x^{n-1}$; el coeficient es multiplica;
        el terme independent es deriva a $0$; terme a terme.
  \item \textbf{Creixement:} $f'(x)>0$ creix, $f'(x)<0$ decreix; tria un valor de prova
        de cada tram.
  \item \textbf{Màxim/mínim:} on $f'(x)=0$; si $f'$ passa de $+$ a $-$ és \textbf{màxim},
        de $-$ a $+$ és \textbf{mínim}.
  \item \textbf{Optimització (5 passos):} objectiu i variable $\to$ funció (amb la
        condició) $\to$ derivar $\to$ classificar $\to$ interpretar.
\end{itemize}
\end{idea}

\noindent\textbf{Pla d'estudi personal.} A partir de les teves graelles d'autoavaluació
(sessions 5 i 7) i del simulacre, anota els \textbf{dos o tres aspectes} que has de
repassar amb més atenció:
\begin{center}
\renewcommand{\arraystretch}{1.5}
\begin{tabular}{c p{10.5cm}}
\toprule
\textbf{\#} & \textbf{Què he de repassar (i com)} \\
\midrule
1 & \rule{9.8cm}{0.4pt} \\
2 & \rule{9.8cm}{0.4pt} \\
3 & \rule{9.8cm}{0.4pt} \\
\bottomrule
\end{tabular}
\end{center}

% ---------------------------------------------------------------------
\fullsolucions{Sessió 15 — simulacre}
\begin{solucions}
\begin{sollist}
  \item \textbf{(Bloc 1)} $f'(x)=3x^2-10x+2$; $f'(2)=12-20+2=-6<0$: en $x=2$ la funció
        \textbf{decreix}.
  \item \textbf{(Bloc 2)} $f'(x)=3x^2-6x-9=3(x-3)(x+1)$, s'anul·la a $x=-1$ i $x=3$.
        Creix a $(-\infty,-1)$ ($f'(-2)=15>0$) i $(3,\infty)$ ($f'(4)=15>0$); decreix a
        $(-1,3)$ ($f'(0)=-9<0$).
  \item \textbf{(Bloc 3)} $f'(x)=3x^2-12=3(x-2)(x+2)$, crítics $x=-2$ i $x=2$. Màxim a
        $(-2,21)$ ($f'$ de $+$ a $-$); mínim a $(2,-11)$ ($f'$ de $-$ a $+$).
  \item \textbf{(Bloc 4)} Costat paral·lel $36-2x$; àrea $A(x)=x(36-2x)=36x-2x^2$;
        $A'(x)=36-4x=0\Rightarrow x=9$ (màxim). Mides: $9$ m i $36-18=18$ m; superfície
        màxima $A(9)=162\ \text{m}^2$. Interpretació: amb la tanca disponible, el recinte
        més gran possible fa $9\times 18$ m.
\end{sollist}
\end{solucions}
